problem:  
Let \(G\) be a simply connected **2‑step nilpotent Lie group** with Lie algebra \(\mathfrak{g} = V \oplus Z\), where \(Z = [\mathfrak{g},\mathfrak{g}]\) is central and \(V\) the horizontal subspace.  
Each left‑invariant metric corresponds (up to \(\mathrm{O}(V)\times\mathrm{O}(Z)\)) to a linear map  
\[
J\colon Z \to \mathfrak{so}(V), \qquad \langle J_z v_1, v_2\rangle = \langle z,[v_1,v_2]\rangle .
\]
Under **normalized Ricci flow**, the induced algebraic evolution \(J(t)\) is the normalized negative gradient flow of the real‑analytic potential  
\[
F(J)=\frac{\|\mathrm{Ric}_J\|^2}{\|J\|^4}.
\]
Hence all trajectories are analytic gradient trajectories on a compactifiable analytic variety.  Convergence to a single critical point (a **2‑step nilsoliton**) is guaranteed only when such a soliton exists, but finer quantitative and structural features of this convergence remain unknown.

**Open problem:**  

1. (**Quantitative convergence for the \(J\)–bracket flow**)  
Let \(J(t)\) evolve by normalized bracket flow within a fixed 2‑step nilpotent isomorphism class that admits a nilsoliton \(J_\infty\).  Determine whether the **Łojasiewicz‑type convergence exponent**
   \[
   \|\!J(t)-J_\infty\!\| \le C\,t^{-\alpha}
   \]
   can be chosen **uniformly** for all such Lie groups of a given pair of dimensions \((\dim V,\dim Z)\), or whether arbitrarily small exponents \(\alpha\) (hence arbitrarily slow asymptotics) can occur as the algebraic type degenerates.

2. (**Degenerate Hessian directions**)  
At a nilsoliton \(J_\infty\), the Hessian of \(F\) defines a quadratic form on the space of equivalence classes of infinitesimal deformations.  
Are there 2‑step nilpotent brackets for which this Hessian has a **non‑isolated zero eigenspace** not induced by automorphisms or scaling—that is, do *flat directions* of the potential truly occur in the moduli space of 2‑step nilpotent brackets?

3. (**Singular orbit stratification**)  
When a 2‑step Lie algebra admits no nilsoliton, describe the possible **limit strata** of normalized flows under natural rescaling of \(V\) versus \(Z\):  
can the algebra degenerate along a one‑parameter algebraic family for which the ratio \(\|J_z\|^2/\|z\|^2\) tends to a nontrivial limiting distribution on \(Z\)?  
Equivalently, is there a canonical limiting nilpotent substructure (possibly of lower step) that captures the collapsed Ricci geometry?

Why is it a "good" problem:  
This reformulation respects the established fact that normalized bracket flow on nilpotent Lie groups is a gradient flow, and it pushes beyond that structure to address what is *not* yet understood: **quantitative rates, degeneracies of the analytic potential, and algebraic stratifications** at the boundary of moduli.  

These questions are intrinsically research‑level because they connect analytic tools (Łojasiewicz inequalities and Hessian spectral theory) with **Lie‑algebraic geometry** (orbit closures and degenerations of \(J\)-maps).  
They bridge **geometric analysis**, **dynamical systems**, and **representation theory**, seeking to reveal the fine structure—rates, stability, and boundary behavior—of Ricci‑flow convergence in the simplest nontrivial nilpotent case.  
The statements are concise and accessible but probe substantial gaps in current knowledge, meeting the criteria of a profound and forward‑looking “good” problem.